% ============================================================
\chapter{Connexion de Levi-Civita}
\label{ch:levi_civita}
% ============================================================

Sur une surface courbe, la notion de \og droite \fg perd son sens euclidien. Mais celle de \og ligne la plus droite possible \fg survit~: c'est la g\'eod\'esique. Pour la d\'efinir, il faut savoir d\'eriver des champs de vecteurs le long de courbes sur la vari\'et\'e --- et c'est l\`a qu'intervient la connexion de Levi-Civita, du nom du math\'ematicien italien Tullio Levi-Civita qui, avec Gregorio Ricci-Curbastro, a d\'evelopp\'e le calcul tensoriel au tournant du XX\textsuperscript{e} si\`ecle. Le r\'esultat fondamental de ce chapitre est aussi surprenant qu'\'el\'egant~: il existe une \emph{unique} fa\c{c}on de d\'eriver qui soit compatible avec la m\'etrique et sans torsion. Cette unicit\'e est le th\'eor\`eme fondamental de la g\'eom\'etrie riemannienne.

\section{Le th\'eor\`eme fondamental de la g\'eom\'etrie riemannienne}

\begin{theoreme}[Th\'eor\`eme fondamental]
\label{thm:levi_civita}
Soit $(M, g)$ une vari\'et\'e riemannienne. Il existe une unique connexion $\nabla$
sur $TM$ v\'erifiant~:
\begin{enumerate}
    \item \textbf{Compatibilit\'e avec la m\'etrique~:} $\nabla g = 0$, i.e.\
    $X\inner{Y}{Z} = \inner{\nabla_X Y}{Z} + \inner{Y}{\nabla_X Z}$.
    \item \textbf{Absence de torsion~:} $\nabla_X Y - \nabla_Y X = [X, Y]$.
\end{enumerate}
Cette connexion est appel\'ee la \emph{connexion de Levi-Civita} de $(M, g)$.
\end{theoreme}

\begin{proof}
\textbf{Unicit\'e.} Supposons que $\nabla$ v\'erifie (1) et (2). En \'ecrivant la
compatibilit\'e pour trois permutations cycliques~:
\begin{align*}
    X\inner{Y}{Z} &= \inner{\nabla_X Y}{Z} + \inner{Y}{\nabla_X Z}, \\
    Y\inner{Z}{X} &= \inner{\nabla_Y Z}{X} + \inner{Z}{\nabla_Y X}, \\
    Z\inner{X}{Y} &= \inner{\nabla_Z X}{Y} + \inner{X}{\nabla_Z Y}.
\end{align*}
En effectuant la combinaison $(1) + (2) - (3)$ et en utilisant l'absence de torsion
($\nabla_X Y - \nabla_Y X = [X,Y]$), on obtient la \emph{formule de Koszul}~:
\begin{align}
    2\inner{\nabla_X Y}{Z} &= X\inner{Y}{Z} + Y\inner{Z}{X} - Z\inner{X}{Y}
    \notag \\
    &\quad + \inner{[X,Y]}{Z} - \inner{[Y,Z]}{X} + \inner{[Z,X]}{Y}.
    \label{eq:koszul}
\end{align}
Le membre de droite est enti\`erement d\'etermin\'e par $g$ et les crochets de Lie,
donc $\nabla_X Y$ est uniquement d\'etermin\'e par la non-d\'eg\'en\'erescence de $g$.

\textbf{Existence.} On d\'efinit $\nabla_X Y$ par la formule de Koszul~\eqref{eq:koszul}.
On v\'erifie que cela d\'efinit bien une connexion (lin\'earit\'e en $X$, r\`egle de Leibniz
en $Y$), qu'elle est compatible avec $g$, et sans torsion. Les v\'erifications sont
directes mais calculatoires.
\end{proof}

\begin{formules}[title=\textbf{Formule de Koszul}]
\[
    2\inner{\nabla_X Y}{Z} = X\inner{Y}{Z} + Y\inner{X}{Z} - Z\inner{X}{Y}
    + \inner{[X,Y]}{Z} - \inner{[X,Z]}{Y} - \inner{[Y,Z]}{X}.
\]
C'est \emph{la} formule fondamentale de la g\'eom\'etrie riemannienne.
\end{formules}

\section{Symboles de Christoffel de la connexion de Levi-Civita}

\begin{proposition}[Formule des symboles de Christoffel]
En coordonn\'ees locales $(x^1, \ldots, x^n)$, les symboles de Christoffel de
la connexion de Levi-Civita sont~:
\[
    \Gamma^k_{ij} = \frac{1}{2} g^{kl}
    \left(\frac{\partial g_{jl}}{\partial x^i}
    + \frac{\partial g_{il}}{\partial x^j}
    - \frac{\partial g_{ij}}{\partial x^l}\right).
\]
\end{proposition}

\begin{proof}
La formule de Koszul appliqu\'ee aux champs coordonn\'ees $\partial_i$, $\partial_j$,
$\partial_l$ (dont les crochets de Lie sont nuls) donne~:
\[
    2\inner{\nabla_{\partial_i}\partial_j}{\partial_l}
    = \partial_i g_{jl} + \partial_j g_{il} - \partial_l g_{ij}.
\]
Donc $2\, g_{kl}\, \Gamma^k_{ij} = \partial_i g_{jl} + \partial_j g_{il} - \partial_l g_{ij}$,
et en multipliant par $g^{lm}/2$, on obtient la formule.
\end{proof}

\begin{remarque}
La sym\'etrie $\Gamma^k_{ij} = \Gamma^k_{ji}$ est imm\'ediate, refl\'etant l'absence
de torsion~: $[\partial_i, \partial_j] = 0$ implique
$\nabla_{\partial_i}\partial_j = \nabla_{\partial_j}\partial_i$.
\end{remarque}

\section{Calculs sur les exemples fondamentaux}

\begin{exemple}[Sph\`ere $S^2$ en coordonn\'ees sph\'eriques]
Avec $g = d\theta^2 + \sin^2\!\theta\, d\phi^2$, on a~:
\[
    g_{11} = 1, \quad g_{22} = \sin^2\!\theta, \quad g_{12} = 0.
\]
Les symboles de Christoffel non nuls sont~:
\begin{align*}
    \Gamma^1_{22} &= -\sin\theta\cos\theta, \\
    \Gamma^2_{12} = \Gamma^2_{21} &= \cot\theta.
\end{align*}
Tous les autres sont nuls. On v\'erifie par la formule~:
$\Gamma^1_{22} = \frac{1}{2}g^{11}(-\partial_1 g_{22}) = -\frac{1}{2}\cdot 2\sin\theta\cos\theta
= -\sin\theta\cos\theta$.
\end{exemple}

\begin{exemple}[Espace hyperbolique $\mathbb{H}^2$ (demi-plan)]
Avec $g = \frac{dx^2 + dy^2}{y^2}$, on a~:
\[
    g_{11} = g_{22} = \frac{1}{y^2}, \quad g_{12} = 0,
    \quad g^{11} = g^{22} = y^2, \quad g^{12} = 0.
\]
Les symboles non nuls sont~:
\begin{align*}
    \Gamma^1_{12} = \Gamma^1_{21} &= -\frac{1}{y}, \\
    \Gamma^2_{11} &= \frac{1}{y}, \\
    \Gamma^2_{22} &= -\frac{1}{y}.
\end{align*}
\end{exemple}

\begin{exemple}[M\'etrique de surface de r\'evolution]
Soit $S$ la surface de r\'evolution param\'etr\'ee par
$(u, v) \mapsto (f(u)\cos v, f(u)\sin v, h(u))$ avec $f > 0$.
La m\'etrique induite est $g = (f'^2 + h'^2)\, du^2 + f^2\, dv^2$.

Si la courbe g\'en\'eratrice est param\'etr\'ee par longueur d'arc ($f'^2 + h'^2 = 1$),
alors $g = du^2 + f(u)^2\, dv^2$ et les symboles non nuls sont~:
\[
    \Gamma^1_{22} = -f f', \qquad
    \Gamma^2_{12} = \Gamma^2_{21} = \frac{f'}{f}.
\]
\end{exemple}

\begin{exemple}[Groupe de Lie avec m\'etrique bi-invariante]
\label{ex:bi_inv_christoffel}
Sur un groupe de Lie $G$ avec m\'etrique bi-invariante, pour deux champs invariants
\`a gauche $X, Y \in \mathfrak{g}$, la formule de Koszul donne~:
\[
    \nabla_X Y = \frac{1}{2}[X, Y].
\]
En effet, la bi-invariance implique $\inner{[X,Y]}{Z} + \inner{Y}{[X,Z]} = 0$
(ad-antisym\'etrie), et les termes $X\inner{Y}{Z}$, etc.\ sont nuls car les champs
sont invariants \`a gauche et la m\'etrique aussi. La formule de Koszul se r\'eduit \`a
$2\inner{\nabla_X Y}{Z} = \inner{[X,Y]}{Z} + \inner{[Z,X]}{Y} - \inner{[Y,Z]}{X}
= \inner{[X,Y]}{Z}$.
\end{exemple}

\section{Gradient, divergence et laplacien}

\begin{definition}[Gradient]
Le \emph{gradient} d'une fonction $f \in C^\infty(M)$ est l'unique champ de
vecteurs $\operatorname{grad} f = (\nabla f)^\sharp$ tel que~:
\[
    g(\operatorname{grad} f, X) = Xf = df(X), \quad \forall\, X \in \mathfrak{X}(M).
\]
En coordonn\'ees~: $(\operatorname{grad} f)^i = g^{ij} \frac{\partial f}{\partial x^j}$.
\end{definition}

\begin{definition}[Divergence]
La \emph{divergence} d'un champ de vecteurs $X \in \mathfrak{X}(M)$ est~:
\[
    \operatorname{div} X = \tr(\nabla X) = \nabla_i X^i
    = \frac{1}{\sqrt{\det g}} \frac{\partial}{\partial x^i}
    \left(\sqrt{\det g}\, X^i\right).
\]
\end{definition}

\begin{definition}[Laplacien de Beltrami]
Le \emph{laplacien} (de Laplace--Beltrami) d'une fonction $f$ est~:
\[
    \Delta f = \operatorname{div}(\operatorname{grad} f)
    = \frac{1}{\sqrt{\det g}} \frac{\partial}{\partial x^i}
    \left(\sqrt{\det g}\, g^{ij} \frac{\partial f}{\partial x^j}\right).
\]
\end{definition}

\begin{exemple}[Laplacien sur $S^2$]
En coordonn\'ees sph\'eriques, $\sqrt{\det g} = \sin\theta$, d'o\`u~:
\[
    \Delta_{S^2} f = \frac{1}{\sin\theta}\frac{\partial}{\partial\theta}
    \left(\sin\theta\, \frac{\partial f}{\partial\theta}\right)
    + \frac{1}{\sin^2\!\theta}\frac{\partial^2 f}{\partial\phi^2}.
\]
\end{exemple}

\begin{exemple}[Laplacien sur $\mathbb{H}^2$]
Avec la m\'etrique $g = y^{-2}(dx^2 + dy^2)$~:
\[
    \Delta_{\mathbb{H}^2} f = y^2\left(\frac{\partial^2 f}{\partial x^2}
    + \frac{\partial^2 f}{\partial y^2}\right).
\]
\end{exemple}

\section{D\'erivation covariante le long d'une courbe}

\begin{definition}
Soit $\gamma : [a,b] \to M$ une courbe lisse et $V$ un champ de vecteurs le long
de $\gamma$. La \emph{d\'eriv\'ee covariante} de $V$ le long de $\gamma$ est~:
\[
    \frac{DV}{dt} = \nabla_{\dot\gamma} V.
\]
En coordonn\'ees~: $\displaystyle\left(\frac{DV}{dt}\right)^k
= \frac{dV^k}{dt} + \Gamma^k_{ij}(\gamma(t))\, \dot\gamma^i(t)\, V^j(t)$.
\end{definition}

\begin{proposition}[R\`egle de Leibniz le long d'une courbe]
Si $V$ et $W$ sont des champs de vecteurs le long de $\gamma$ et $\nabla$ est
compatible avec $g$, alors~:
\[
    \frac{d}{dt}\inner{V}{W} = \inner{\frac{DV}{dt}}{W} + \inner{V}{\frac{DW}{dt}}.
\]
\end{proposition}

\section{Acc\'el\'eration covariante et g\'eod\'esiques}

\begin{definition}[Acc\'el\'eration covariante]
L'\emph{acc\'el\'eration covariante} d'une courbe $\gamma$ est~:
\[
    \frac{D\dot\gamma}{dt} = \nabla_{\dot\gamma}\dot\gamma.
\]
En coordonn\'ees~: $\displaystyle\left(\frac{D\dot\gamma}{dt}\right)^k
= \ddot\gamma^k + \Gamma^k_{ij}\, \dot\gamma^i\, \dot\gamma^j$.
\end{definition}

\begin{definition}[G\'eod\'esique]
Une courbe $\gamma$ est une \emph{g\'eod\'esique} si son acc\'el\'eration covariante
est nulle~:
\[
    \frac{D\dot\gamma}{dt} = \nabla_{\dot\gamma}\dot\gamma = 0.
\]
\end{definition}

L'\'equation des g\'eod\'esiques en coordonn\'ees est le syst\`eme d'EDO du second ordre~:
\[
    \ddot\gamma^k + \Gamma^k_{ij}(\gamma)\, \dot\gamma^i\, \dot\gamma^j = 0,
    \quad k = 1, \ldots, n.
\]

\begin{formules}[title=\textbf{\'Equation des g\'eod\'esiques}]
\[
    \ddot\gamma^k + \Gamma^k_{ij}\, \dot\gamma^i\, \dot\gamma^j = 0.
\]
C'est un syst\`eme de $n$ EDO du second ordre, \'equivalent \`a $2n$ EDO du premier
ordre sur $TM$.
\end{formules}

\section{Connexion de Levi-Civita des produits et produits tordus}

\begin{proposition}[Produit riemannien]
Sur $(M_1 \times M_2, g_1 \oplus g_2)$, si $X_1, Y_1 \in \mathfrak{X}(M_1)$ et
$X_2, Y_2 \in \mathfrak{X}(M_2)$, alors~:
\[
    \nabla_{(X_1, X_2)}(Y_1, Y_2) = (\nabla^1_{X_1} Y_1, \nabla^2_{X_2} Y_2),
\]
o\`u $\nabla^1$ et $\nabla^2$ sont les connexions de Levi-Civita de $(M_1, g_1)$
et $(M_2, g_2)$ respectivement.
\end{proposition}

\begin{proposition}[Produit tordu]
\label{prop:warped_LC}
Sur le produit tordu $B \times_f F$ avec $g = g_B + f^2 g_F$, la connexion de
Levi-Civita v\'erifie, pour $X, Y \in \mathfrak{X}(B)$ et $V, W \in \mathfrak{X}(F)$~:
\begin{align*}
    \nabla_X Y &= \nabla^B_X Y, \\
    \nabla_X V = \nabla_V X &= \frac{Xf}{f}\, V, \\
    \nabla_V W &= \nabla^F_V W - \frac{\inner{V}{W}_F}{f}\, \operatorname{grad}_B f.
\end{align*}
\end{proposition}

\section{Exercices}

\begin{exercice}
Calculer tous les symboles de Christoffel pour la m\'etrique
$g = dr^2 + r^2(d\theta^2 + \sin^2\!\theta\, d\phi^2)$ de $\R^3$ en coordonn\'ees
sph\'eriques.
\end{exercice}

\begin{exercice}
V\'erifier que les grands cercles de $S^2$ sont bien des g\'eod\'esiques en utilisant
les symboles de Christoffel calcul\'es ci-dessus.
\end{exercice}

\begin{exercice}
Montrer que sur un groupe de Lie avec m\'etrique bi-invariante, les g\'eod\'esiques
passant par l'\'el\'ement neutre $e$ sont exactement les sous-groupes \`a un param\`etre
$t \mapsto \exp(tX)$.
\end{exercice}

\begin{exercice}
Calculer le laplacien de la fonction $f(x,y) = \ln y$ sur $\mathbb{H}^2$ muni
de la m\'etrique $g = \frac{dx^2 + dy^2}{y^2}$.
\end{exercice}

\begin{exercice}
Soit $(M, g)$ riemannienne et $\tilde{g} = e^{2\varphi} g$ un changement conforme.
Montrer que les symboles de Christoffel sont reli\'es par~:
\[
    \tilde{\Gamma}^k_{ij} = \Gamma^k_{ij} + \delta^k_i \partial_j \varphi
    + \delta^k_j \partial_i \varphi - g_{ij} g^{kl} \partial_l \varphi.
\]
\end{exercice}

\begin{exercice}
Sur le produit tordu $\R \times_{\cosh t} \R^{n-1}$ muni de
$g = dt^2 + \cosh^2(t)\, g_{\R^{n-1}}$, calculer les symboles de Christoffel
et v\'erifier que les courbes $t \mapsto (t, x_0)$ sont des g\'eod\'esiques.
\end{exercice}
